% ==============================================================
\section{関連研究}\label{sec:related}
% ==============================================================

\subsection{頻出パターンマイニング}

頻出パターンマイニングは，トランザクションデータの一定割合以上に出現するアイテムセットを発見する技術である．
Apriori~\cite{agrawal1994fast}は反単調性に基づく候補生成・枝刈りフレームワークを確立した．
FP-Growth~\cite{han2000mining}はFP-treeにより候補生成を回避し高速化を実現した．
ECLAT~\cite{zaki2000scalable}は縦型（tidset）表現と集合積により効率的なサポート計算を行う．
LCM~\cite{uno2004lcm}は頻出・閉・極大アイテムセットの最適列挙を達成した．
これらの手法はグローバル頻度を対象とし，時間的な変動をモデル化しない．

\subsection{時間的・局所的パターンマイニング}

頻出パターンの時間的側面を扱う拡張研究がいくつか存在する．
LPFIM~\cite{mahanta2005finding}は時間区間内の局所頻出集合を対象とするが，
gap ベースの閾値に依存する．
PFPM~\cite{tanbeer2009discovering}およびPPFPM~\cite{kiran2017discovering}は
周期性制約を導入した．
LPPM~\cite{fournier2021mining}は局所性と周期性の基準を統合する．
Apriori-window~\cite{dsaa2025}は固定幅スライディングウィンドウ上で
局所サポートを直接評価し，gapパラメータに依存しない密集パターン抽出を実現した．

これらの手法はいずれもパターンが\emph{いつ}頻出であるかを検出するが，
頻度が変動する\emph{原因}はモデル化しない．
本研究はこれらを補完し，検出されたサポート変動を外部イベントに帰属させる．

\subsection{変化点検出}

変化点検出は時系列データの構造変化を同定する技術である．
CUSUMアルゴリズム~\cite{page1954continuous}は逐次過程の平均シフトを検出する．
包括的なサーベイ~\cite{aminikhanghahi2017survey,truong2020selective}が
パラメトリック・ノンパラメトリック・オンライン手法を網羅している．
これらの手法は個別時系列に適用されるが，
数千のパターンサポート時系列を同時分析する多重検定の負荷には対応せず，
検出された変化をイベントに結びつける機能も持たない．

\subsection{イベント帰属と因果推論}

計量経済学のイベント・スタディ法~\cite{mackinlay1997event}は，
企業イベント（合併，決算発表等）の株価への影響を
異常リターン分析により定量化する．
分散システムにおけるroot cause analysis~\cite{chen2004pinpoint,jeyakumar2019explainit}は，
異常を特定のコンポーネントや設定変更に帰属させる．
Granger因果~\cite{granger1969investigating}とその拡張は，
ある時系列が別の時系列の予測に寄与するか否かを検定する．

これらのアプローチは典型的には単一または少数のアウトカム変数を分析する．
対照的に，本研究の設定では数千のパターンサポート時系列を複数のイベントに対して
同時に検定する必要があり，偽発見率を制御するための多重検定補正が不可欠である．

\subsection{置換検定と多重検定補正}

置換検定~\cite{good2005permutation}は，
データの並べ替えから帰無分布を構築する分布自由な有意性検定法である．
複数仮説の同時検定においては，
Bonferroni補正~\cite{dunn1961multiple}が族単位誤り率（FWER）を制御するが保守的であり，
Benjamini--Hochberg (BH) 手続き~\cite{benjamini1995controlling}が
偽発見率（FDR）を制御しつつより高い検出力を提供する．

本パイプラインは，円形シフト置換検定~\cite{davison1997bootstrap}と
全パターン・全イベントにわたるグローバルBH補正を組み合わせ，
分布自由な検定と適切な多重検定制御を両立させる．
